\newcommand{\scribefigureplaceholder}[2]{%
\begin{figure}[ht]
  \centering
  \fbox{%
    \parbox[c][2.3in][c]{0.84\textwidth}{%
      \centering
      \textbf{Illustration Placeholder}\\[0.5em]
      #2
    }%
  }
  \caption{#1}
\end{figure}
}

\newcommand{\scribecodeplaceholder}[2]{%
\begin{ppprogram}{#1}{#2}
  \centering
  \fbox{%
    \parbox[c][2.5in][c]{0.90\linewidth}{%
      \centering
      \textbf{Code Placeholder}\\[0.5em]
      Insert the final code or pseudocode from the lecture demo here.
    }%
  }
\end{ppprogram}
}

\chapter{Performance Metrics and Basic PRAM Algorithms}

\section{Metrics We Will Use}
For an input of size $n$, let $T(n,1)$ denote the best sequential running time
and $T(n,p)$ the running time on $p$ processors. The central performance measures
used throughout the lecture are:
\begin{align*}
S(n,p) &= \frac{T(n,1)}{T(n,p)}, \\
C(n,p) &= p \cdot T(n,p), \\
E(n,p) &= \frac{S(n,p)}{p} = \frac{T(n,1)}{C(n,p)}.
\end{align*}
Here $S(n,p)$ is the speedup, $C(n,p)$ is the cost, and $E(n,p)$ is the efficiency.
The parallelism overhead is
\[
T_o(n,p) = C(n,p) - T(n,1).
\]
An algorithm is \ppterm{cost-optimal} when $C(n,p) = O(T(n,1))$, and is
\ppterm{work-efficient} when the total number of primitive operations remains
within a constant factor of the best sequential algorithm. These notions will
be used repeatedly when comparing search, reduction, and prefix-sum algorithms
\citep{amdahl1967,gustafson1988,pacheco2011}.

\begin{ppkey}
Fixed-size speedup is limited by the sequential fraction of the computation,
but scalability improves when the problem size grows with the number of processors.
\end{ppkey}

\section{The PRAM Abstraction}
The lecture uses the PRAM model as an idealised machine for algorithm design:
processors are synchronous, each processor has a private identity and local
memory, and shared-memory accesses are treated as constant-time operations.
Although unrealistic as a hardware model, PRAM is extremely useful for asymptotic
reasoning.

\begin{ppdefn}
\ppterm{EREW} (exclusive read, exclusive write) allows at most one processor to
read or write a shared location in a step. \ppterm{CREW} permits concurrent
reads but still requires exclusive writes. \ppterm{CRCW} permits both concurrent
reads and concurrent writes, but a write-conflict rule must then be specified.
\end{ppdefn}

Under CRCW, common variants include \ppterm{priority write}, where the highest-priority
processor wins; \ppterm{common write}, where all concurrent writers must write the
same value; and \ppterm{arbitrary write}, where one of the proposed values is chosen
arbitrarily. The choice of model matters because it changes which synchronisation
patterns are legal and how much coordination is needed.

\section{Parallel Linear Search}
Suppose an unsorted array $A[1 \dots n]$ is searched for a value $x$ using $p$
processors. A natural block-partitioned strategy assigns roughly $n/p$ elements
to each processor. Each processor performs a local linear scan on its own block
and produces one Boolean answer: whether $x$ was found in that block.

The local scan takes $\Theta(n/p)$ time. These local answers can then be combined
using a tree reduction with logical OR, which takes $\Theta(\log p)$ time. Hence
\[
T(n,p) = \Theta\left(\frac{n}{p} + \log p\right).
\]
The corresponding cost is
\[
C(n,p) = p \cdot T(n,p) = \Theta(n + p \log p).
\]
Therefore the algorithm is cost-optimal when $p \log p = O(n)$; in the common
regime $n \gg p$, the scan term dominates and this reduction-based solution is
close to work-efficient.

\begin{ppnote}
The tree-reduction version is EREW-safe because processors read disjoint input
blocks and later reduce flags through disjoint pairs. No concurrent access to
the same shared location is required.
\end{ppnote}

The lecture also discussed a stronger shared-memory variant in which processors
write to a shared Boolean register such as \ppinline{found}. In a common-CRCW
model, all successful processors may write the value \ppinline{true}, so the final
combine can be reduced to constant time after the scan phase. This is faster,
but it relies on a stronger concurrency model than the reduction-based EREW version.

An alternate EREW-friendly strategy is to give one element to each processor and
write hits into a private \ppinline{hit[i]} location. A subsequent tree OR-reduction
over the hit array decides whether any processor found the target. With $n$ processors,
this gives $\Theta(\log n)$ time and cleanly separates the private-write phase from
the reduction phase.

\scribefigureplaceholder{Block-partitioned search over an array}{
Show an array divided into $p$ contiguous blocks, one processor per block, and a
tree OR-reduction over the local Boolean answers.
}

\scribecodeplaceholder{Block-partitioned linear search with OR-reduction}{prog:block-search-placeholder}

\section{Tree-Based Parallel Addition}
A naive parallelisation of summation would let all processors update one shared
\ppinline{sum} variable. This is a poor design: the updates must be serialised or
protected, so the resulting algorithm gains little or no speedup.

The correct approach is a reduction tree. Assume for simplicity that $n$ is a power
of two. In round 1, adjacent pairs are added. In round 2, the stride doubles and
the partial sums of length two are combined into sums of length four. After
$\log_2 n$ rounds, one processor holds the total sum of all elements.

With one processor per active array position in each round, the parallel time is
$\Theta(\log n)$. The total amount of arithmetic work is
\[
\frac{n}{2} + \frac{n}{4} + \cdots + 1 = \Theta(n),
\]
which matches sequential summation up to constant factors. However, if $n$ processors
are reserved throughout all $\log n$ rounds, the cost becomes $\Theta(n \log n)$,
so the algorithm is not cost-optimal in that raw form even though the arithmetic
work itself is linear.

For $p < n$, a standard cost-optimal version first computes local sums over blocks
of size about $n/p$, and then reduces the $p$ block sums in $\Theta(\log p)$ more
steps. This yields
\[
T(n,p) = \Theta\left(\frac{n}{p} + \log p\right),
\]
which matches the block-partitioned search bound.

\begin{pppitfall}
Work and cost are not the same quantity. A reduction tree performs only $\Theta(n)$
additions, but if too many processors sit idle across too many rounds, the total
cost can still be $\Theta(n \log n)$.
\end{pppitfall}

\scribefigureplaceholder{Reduction tree for parallel addition}{
Illustrate pairwise additions across rounds, with the stride doubling at each level
until a single root stores the total sum.
}

\scribecodeplaceholder{Stride-based tree reduction for summation}{prog:parallel-sum-placeholder}

\section{Parallel Prefix Sums}
The prefix-sum or scan problem asks us to transform an input array
$A[1 \dots n]$ into prefix sums. In the \ppterm{inclusive} version,
\[
B[i] = \sum_{j=1}^{i} A[j],
\]
whereas in the \ppterm{exclusive} version,
\[
B[i] = \sum_{j=1}^{i-1} A[j],
\]
with the convention that $B[1] = 0$. Sequentially, either version takes $\Theta(n)$
time.

\subsection{Hillis-Steele Scan}
The Hillis-Steele scan proceeds in rounds. At round $d$, processor $i$ adds into
its current value the entry from position $i - 2^{d-1}$, provided that this index
exists. Conceptually, each round must read values from the \emph{previous} round,
not partially updated values from the current round.

\paragraph{Correctness proof by induction.}
Let $B^{(0)}[i] = A[i]$. For each round $d \ge 1$, define
\[
B^{(d)}[i] =
\begin{cases}
B^{(d-1)}[i] + B^{(d-1)}[i - 2^{d-1}], & \text{if } i > 2^{d-1}, \\
B^{(d-1)}[i], & \text{otherwise}.
\end{cases}
\]
We claim that after round $d$, every index $i$ stores the sum of the last
$\min(2^d,i)$ input values ending at position $i$:
\[
B^{(d)}[i] = \sum_{j=\max(1,\, i - 2^d + 1)}^{i} A[j].
\]
The proof is by induction on $d$.

For the base case $d = 0$, we have $B^{(0)}[i] = A[i]$, which is exactly the sum
over the last $\min(1,i) = 1$ entry ending at $i$.

Now assume the claim holds after round $d-1$. Fix an index $i$.
If $i \le 2^{d-1}$, then no update happens in round $d$, so
$B^{(d)}[i] = B^{(d-1)}[i]$. By the induction hypothesis this is already
\[
\sum_{j=1}^{i} A[j]
= \sum_{j=\max(1,\, i - 2^d + 1)}^{i} A[j],
\]
because $i - 2^d + 1 \le 1$ in this case.

If instead $i > 2^{d-1}$, then by the update rule and the induction hypothesis,
\begin{align*}
B^{(d)}[i]
&= B^{(d-1)}[i] + B^{(d-1)}[i - 2^{d-1}] \\
&= \sum_{j=i-2^{d-1}+1}^{i} A[j]
 + \sum_{j=\max(1,\, i-2^d+1)}^{i-2^{d-1}} A[j] \\
&= \sum_{j=\max(1,\, i-2^d+1)}^{i} A[j].
\end{align*}
The two sums are over adjacent ranges, so together they cover exactly the last
$\min(2^d,i)$ entries ending at position $i$.

Therefore the claim holds for every round $d$. After
$D = \lceil \log_2 n \rceil$ rounds we have $2^D \ge n$, so for every $i$,
\[
B^{(D)}[i] = \sum_{j=1}^{i} A[j],
\]
which is precisely the inclusive prefix sum. Hence the Hillis-Steele algorithm
is correct.

For the example array
\[
[3,1,4,5,1,9,2],
\]
the inclusive scan evolves as follows:
\[
[3,4,5,9,6,10,11],
\]
\[
[3,4,8,13,11,19,17],
\]
\[
[3,4,8,13,14,23,25].
\]
After $\lceil \log_2 n \rceil$ rounds, every position holds the correct prefix sum.

The time is $\Theta(\log n)$ on $n$ processors, but every round touches $\Theta(n)$
entries, so the total work is $\Theta(n \log n)$. Thus Hillis-Steele is easy to
parallelise, but it is not work-efficient \citep{hillis1986,harris2007scan}.

\scribefigureplaceholder{Round-by-round view of Hillis-Steele scan}{
Show the array after each doubling step $d = 1,2,3$ so that the dependency pattern
from index $i - 2^{d-1}$ is visually clear.
}

\scribecodeplaceholder{Hillis-Steele inclusive scan}{prog:hillis-steele-placeholder}

\subsection{Blelloch Scan}
Blelloch's algorithm improves the work bound by using two tree phases
\citep{blelloch1990,harris2007scan}. For simplicity, assume the input size is a
power of two; otherwise pad the array with zeros up to the next power of two.

In the \ppterm{upsweep} (or reduction) phase, a balanced binary tree is built over
the array. Each internal node stores the sum of its two children, so the root
eventually stores the total sum of all elements. This phase resembles tree-based
parallel addition.

In the \ppterm{downsweep} phase, the root is first set to $0$. Then each internal
node propagates prefix information downward as follows. If the parent currently
stores $p$ and the left child previously stored a subtotal $t$, then the left child
receives $p$ and the right child receives $p + t$. Repeating this from the root
to the leaves yields the \emph{exclusive} scan.

The inclusive scan can be recovered by adding the original input element-wise to
the exclusive result. Across both phases, the total work is $\Theta(n)$ and the
parallel time is $\Theta(\log n)$ on $n$ processors, so Blelloch scan is work-efficient.

\begin{ppkey}
Hillis-Steele gives a simple scan in $\Theta(\log n)$ time but spends $\Theta(n \log n)$
work. Blelloch keeps the same asymptotic time while reducing the work to $\Theta(n)$.
\end{ppkey}

\scribefigureplaceholder{Blelloch upsweep and downsweep}{
Draw a binary tree over the array, mark subtree sums during the upsweep, then show
how setting the root to $0$ and propagating offsets downward produces the exclusive scan.
}

\scribecodeplaceholder{Blelloch exclusive scan with an inclusive conversion step}{prog:blelloch-placeholder}

\section{Summary}
Three themes recur across all of these examples. First, the computation model
matters: the legality and cost of synchronisation depend on whether we are in EREW,
CREW, or CRCW. Second, speedup alone is not enough; cost and work-efficiency are
equally important. Third, good parallel algorithms combine local work with a small
number of structured synchronisation phases, typically reductions or scans.
